%% kapa: exam=1 round=29 year=2018 period=1 form=A subject=economics subject_ko=경제학원론 seq=2 number=42
%% source: Q-Net 감정평가사 기출문제 (cst00309) artlSeq=5202808 file=extracted/1008/5202808/2018년 제29회 감정평가사 1차시험 기출문제/2018년 제29회 감정평가사 1차시험 기출문제/1교시 A형.pdf
%% conversion: kordoc 4.24.51 -> tools/gen1.py (auto); review=proofread-sonnet; flags=none
\begin{question}{42}
밑줄 친 변화에 따라 각국의 노동시장에서 예상되는 현상으로 옳은 것은? (단, 노동수요곡선은 우하향, 노동공급곡선은 우상향하고, 다른 조건은 일정하다.)
\begin{framedtext}
○ 甲국에서는 (A)\uline{인구 감소로 노동시장에 참여하고자 하는 사람들이 감소}하였다. \par
○ 乙국의 정부는 (B)\uline{규제가 없는 노동시장에 균형임금보다 높은 수준에서 최저임금제를 도입}하려고 한다. \par
\end{framedtext}
\begin{choices}
  \item (A): 노동수요 감소, (B): 초과수요 발생
  \item (A): 노동수요 증가, (B): 초과공급 발생
  \item (A): 노동공급 감소, (B): 초과수요 발생
  \item (A): 노동공급 증가, (B): 초과공급 발생
  \item (A): 노동공급 감소, (B): 초과공급 발생
\end{choices}
\end{question}
\answer{5}
